
Die bisherige Betrachtung verdeutlicht, dass ERP-Systeme zwar ein hohes Potenzial für die digitale Transformation von 
KMU besitzen, ihre Wirkung jedoch nicht automatisch aus der reinen Einführung der Software entsteht. Ein ERP-System 
kann Prozesse integrieren, Daten zentralisieren und Transparenz schaffen. Ob diese Potenziale tatsächlich realisiert 
werden, hängt jedoch wesentlich davon ab, ob bestimmte organisatorische, technologische und strategische Voraussetzungen 
erfüllt sind. Die Wirksamkeit eines ERP-Systems ist somit nicht allein eine Frage der technischen Funktionalität, 
sondern insbesondere eine Frage der Passung zwischen System, Organisation und Unternehmenszielen \parencite{Leyh2015,Alaskari2021,Leyh2019}.

Eine erste zentrale Voraussetzung besteht darin, dass die Geschäftsprozesse eines Unternehmens ERP-fähig sind. 
ERP-Systeme beruhen in hohem Maße auf der Standardisierung und Integration betrieblicher Abläufe. Damit ein solches 
System seine Wirkung entfalten kann, müssen die relevanten Prozesse ausreichend klar beschrieben, strukturiert und 
im System abbildbar sein \parencite{Leyh2019}. Sind Prozesse dagegen stark informell, uneinheitlich oder nur schwer standardisierbar, 
steigt häufig der Bedarf an individuellen Anpassungen. Solche Anpassungen können insbesondere für KMU problematisch 
sein, da sie zusätzliche Kosten, technische Komplexität und langfristige Abhängigkeiten verursachen können 
\parencite{Leyh2015}. Prozessklarheit bildet daher eine wichtige Grundlage für die erfolgreiche Nutzung eines ERP-Systems.

Eng damit verbunden ist die Qualität und Nutzbarkeit der zugrunde liegenden Daten. ERP-Systeme entfalten ihren Nutzen 
vor allem dadurch, dass sie bisher getrennte Informationsbestände zusammenführen und eine gemeinsame Datenbasis 
schaffen \parencite{Tereshchenko2016}. Diese zentrale Datenbasis kann Transparenz erhöhen, Abstimmungsaufwand reduzieren und datenbasierte 
Entscheidungen unterstützen. Voraussetzung dafür ist jedoch, dass die im System verarbeiteten Daten konsistent, 
aktuell und zuverlässig sind. Inkonsistente Stammdaten, doppelte Datenpflege oder fehlerhafte Datenmigration 
können die Integrationswirkung eines ERP-Systems erheblich einschränken \parencite{Leyh2015,Leyh2019}. Die 
Datenqualität ist damit ein zentraler Bestandteil der ERP-Wirksamkeit.

Darüber hinaus müssen die für ein ERP-Projekt erforderlichen Ressourcen realistisch verfügbar sein. ERP-Einführungen 
erfordern nicht nur finanzielle Mittel, sondern auch Zeit, personelle Kapazitäten und die aktive Beteiligung 
verschiedener Fachbereiche \parencite{Alaskari2021}. Gerade in KMU können begrenzte Budgets und knappe interne Ressourcen den Projekterfolg 
gefährden \parencite{Leyh2015}. Wird ein ERP-Projekt zu knapp geplant oder unterfinanziert, kann dies andere Erfolgsfaktoren negativ 
beeinflussen, etwa Schulungen, Prozessanalysen oder die Qualität der Systemeinführung \parencite{Xia2009,Barbieri2025}. 
Die Ressourcenverfügbarkeit ist daher nicht nur eine operative Rahmenbedingung, sondern beeinflusst die gesamte Umsetzung 
und spätere Nutzung des Systems.

Ein weiterer Erfolgsfaktor liegt in den Kompetenzen und der Akzeptanz der Mitarbeitenden. ERP-Systeme verändern 
bestehende Arbeitsabläufe und verlangen von den Anwenderinnen und Anwendern ein Verständnis für neue Prozesse, 
Rollen und Informationsflüsse \parencite{Alaskari2021,Leyh2015}. Wenn Mitarbeitende das System nicht verstehen oder dessen Nutzen nicht erkennen, 
besteht die Gefahr, dass Funktionen nicht genutzt, Prozesse umgangen oder alte Arbeitsweisen beibehalten werden \parencite{Leyh2015,Leyh2019}. 
Schulungen, Beteiligung der Fachbereiche und ein aktives Change Management sind daher wesentlich, um Akzeptanz 
aufzubauen und die tatsächliche Nutzung des Systems sicherzustellen \parencite{Noudoostbeni2009,Faruque2024}.

Schließlich ist entscheidend, dass ERP nicht ausschließlich als technisches IT-Projekt verstanden wird. Ein 
ERP-System kann insbesondere dann transformativ wirken, wenn es in eine übergeordnete Unternehmens- und 
Digitalisierungsstrategie eingebettet ist \parencite{Leyh2019}. Dazu gehört, dass die Ziele der ERP-Einführung klar definiert sind 
und mit langfristigen Unternehmenszielen verknüpft werden. Die Unterstützung durch das Management spielt hierbei 
eine zentrale Rolle, da ERP-Projekte häufig bereichsübergreifende Veränderungen auslösen und strategische 
Entscheidungen erfordern \parencite{Leyh2015, Barbieri2025}. Erst wenn ERP als Grundlage weiterer Digitalisierung verstanden 
wird, kann das System über die reine Prozessunterstützung hinaus eine transformierende Wirkung entfalten.

Zusammenfassend zeigt sich, dass die Wirksamkeit von ERP-Systemen in KMU aus dem Zusammenspiel mehrerer 
Voraussetzungen entsteht. Prozesse müssen ERP-fähig sein, Daten müssen integriert und nutzbar vorliegen, 
Ressourcen müssen realistisch geplant werden, Mitarbeitende müssen über Kompetenzen und Akzeptanz verfügen 
und das ERP-System muss strategisch eingebettet sein. 
%Diese Erfolgslogik bildet den Ausgangspunkt für das im folgenden Kapitel entwickelte ERP-KMU-Fit-Modell.
